\chapter{帰納的関数と計算可能関数}
\label[chapter]{chap:computable}

これまでの議論では，数学における「論理」に関連した概念を形式体系として定式化し，その性質を調べてきた．
特に\cref{chap:proof}では，数学における「証明」を形式的な記号列の変形操作として定式化した．
この形式的な記号列の変形操作は「与えられた記号列に対していくつかの変形規則を適用し，目的の記号列に変形する」という形で行われるが，
これは我々が普段行っている「計算」に非常によく似ている．
この関連性は，不完全性定理をめぐる議論においても重要な役割を果たす．

この章では，この「計算ができること」を定式化し，その性質を調べることを試みる．
上記の「計算ができること」の数学的な定式化はいくつも知られている．
本書では，廣瀬\cite{hirose2024}の流れに沿っての2種の定式化を取り扱う．
他にもチューリングマシンによる定式化やラムダ記法による定式化なども知られており，これらの定式化はすべて等価であることが知られている．
「計算ができること」のいくつもの定式化がすべて等価であるという事実は，これらの定式化が「計算ができること」の本質をよくとらえていることを示唆している．
このことから，これらの関数を「計算可能関数」とみなすにしようという%
\index[widx]{Churchのテーゼ}%
Churchのテーゼ%
\footnote{%
	「テーゼ」とは，ものすごく大雑把にいえばある事柄に対する肯定的な主張のことを指す．
}%
がなされ，現代では広く受け入れられている．

\section{帰納的関数}
\label[section]{sec:inductive-function}

まずは，計算可能性の1つ目の定式化として帰納的関数を定義する．

\index[widx]{きのうてきかんすう@帰納的関数}
\index[widx]{きのうてきかんすう@帰納的関数!げんしきのうてきかんすう@原始---}
\index[widx]{ぜろかんすう@ゼロ関数}
\index[widx]{こうしゃかんすう@後者関数}
\index[widx]{しゃえいかんすう@射影関数}
\index[sidx]{\(\muoperator\)：\(\muoperator\)作用素}
\index[widx]{\(\muoperator\)さようそ@\(\muoperator\)さようそ}
\index[sidx]{\(\zero^n\)：ゼロ関数}
\index[sidx]{\(\proj_i^n\)：射影関数}
\index[sidx]{\(\suc\)：後者関数}
\begin{Def} \label{Def:recursivefunction}
	\term{帰納的関数}は以下のように帰納的に定義される：
	\begin{enumerate}
		\item 正の整数\(n\)に対して\(\NaturalNumbers^n\)から\(\NaturalNumbers\)への関数\(x \mapsto 0\)は帰納的関数である．
		      この関数を
		      \begin{equation}
			      \zero^n
			      \label{eq:zerofunction}
		      \end{equation}
		      と表す．
		      各\(n\)に対する\(\zero^n\)を総称してゼロ関数呼ぶ．
		\item \(\NaturalNumbers\)から\(\NaturalNumbers\)への関数\(x \mapsto x + 1\)は帰納的関数である．この関数を後者関数と呼び，
		      \begin{equation}
			      \suc
			      \label{eq:sucfunction}
		      \end{equation}
		      と表す．
		\item \(n\)を正の整数とする．\(i=1,2,\dots, n\)に対して，\(\NaturalNumbers^n\)から\(\NaturalNumbers\)への関数\(\pair{x_1, \dots, x_n} \mapsto x_i\)は帰納的関数である．
		      この関数を
		      \begin{equation}
			      \proj_i^n
		      \end{equation}
		      と表す．
		      各\(i, n\)に対する\(\proj_i^n\)を総称して射影関数と呼ぶ．
		\item \(n,m\)を正の整数とする．\(h \colon \NaturalNumbers^m \to \NaturalNumbers\)が帰納的関数で，
		      \(i = 1,2,\dots,m\)に対して\(g_i \colon \NaturalNumbers^n \to \NaturalNumbers\)が帰納的関数ならば，
		      \begin{align*}
			      \apply{f}{x_1, \dots, x_n} = \apply{h}{
				      \apply{g_1}{x_1, \dots, x_n},
				      \apply{g_2}{x_1, \dots, x_n},
				      \dots,
				      \apply{g_m}{x_1, \dots, x_n}
			      } \\
			      \text{(\(x_1, \dots, x_n \in \NaturalNumbers\))}
		      \end{align*}
		      で定義される関数\(f \colon \NaturalNumbers^n \to \NaturalNumbers\)は帰納的関数である．
		      \(f\)は\(g, h\)から関数合成によって得られる関数である．
		\item \(n\)を正の整数とする．\(g \colon \NaturalNumbers^n \to \NaturalNumbers\)と\(h \colon \NaturalNumbers^{n+2}\)がともに帰納的関数ならば，
		      \begin{align*}
			      \apply{f}{x_1, x_2, \dots, x_n}        & = \apply{g}{x_1, x_2, \dots, x_n} \quad \text{(\(x_1, x_2, \dots, x_n \in \NaturalNumbers\))}, \\
			      \apply{f}{x_1, x_2, \dots, x_n, y + 1} & = \apply{h}{x_1, x_2, \dots, x_n, y, \apply{f}{x_1, x_2, \dots, x_n}}                          \\
			                                             & \text{(\(x_1, x_2, \dots, x_n \in \NaturalNumbers\))}
		      \end{align*}
		      で定義される関数\(f \colon \NaturalNumbers^{n+1} \to \NaturalNumbers\)は帰納的関数である．
		      \(f\)は\(g, h\)から原始帰納法によって定義される関数である．
		\item \(g \colon \NaturalNumbers^{n + 1} \to \NaturalNumbers\)は帰納的関数で，すべての\(x_1, x_2, \dots, x_n\)に対して
		      \(\apply{g}{x_1, x_2, \dots, x_n, y} = 0\)となる\(y\)が存在したとする．
		      そのような\(y\)のうち最小のものを
		      \begin{equation}
			      \apply{\muoperator y}{\apply{g}{x_1, x_2, \dots, x_n, y} = 0}
			      \label{eq:muoperator}
		      \end{equation}
		      と表すこととする．記号\(\muoperator\)を
		      \term{\(\muoperator\)作用素}と呼ぶことがある．
		      このとき，
		      \[
			      \apply{f}{x_1, \dots, x_n} = \apply{\muoperator y}{\apply{g}{x_1, x_2, \dots, x_n, y} = 0}
			      \quad \text{(\(x_1, x_2, \dots, x_n \in \NaturalNumbers\))}
		      \]
		      で定義される関数\(f \colon \NaturalNumbers^{n+1} \to \NaturalNumbers\)は帰納的関数である．
		\item 以上の規則を有限回適用して得られる関数のみが帰納的関数である．
	\end{enumerate}

	また，規則6を1度も適用せずに得られる帰納的関数を特に
	\term{原始帰納的関数}と呼ぶ．
\end{Def}

\begin{Note}
	与えられた帰納的関数について，その構成過程は一意に定まらない．
	例えば，ゼロ関数は何回合成してもゼロ関数である．
	よって，帰納的関数同士の関係としてその定義に付随する順序が整礎であるとはいえない．
	しかし，それぞれの定義同士の関係としての順序は整礎である．
	そのため，帰納的関数の構成に関する帰納法を利用する場合，代わりにその定義の構成に関する帰納法を用いる．
	帰納的関数に関する定理の主張の証明としてはこれで問題ないので，
	帰納的関数の定義の構成に関する帰納法も帰納的関数の構成に関する帰納法と呼ぶことが多い．
\end{Note}

議論の本筋とは少々ずれるが，帰納的関数の定義に関する帰納法の題材として
以下のAckermann関数を考えよう．

\index[widx]{Ackermannかんすう@Ackermann関数}
\index[sidx]{\(\Ack\)：Ackermann関数}
\begin{Def} \label{Def:ackermannfunction}
	Ackermann関数\(\Ack \colon \NaturalNumbers^2 \to \NaturalNumbers\)を以下のように帰納的に定義する：
	\begin{align}
		\begin{aligned}
			\apply{\Ack}{0, y}         & = y + 1 \quad \text{(\(y \in \NaturalNumbers\))},              \\
			\apply{\Ack}{x + 1, 0}     & = \apply{\Ack}{x, 1} \quad \text{(\(x \in \NaturalNumbers\))}, \\
			\apply{\Ack}{x + 1, y + 1} & = \apply{\Ack}{x, \apply{\Ack}{x + 1, y}} \quad
			\text{(\(x, y \in \NaturalNumbers\))}.
		\end{aligned}
		\label{eq:ackermannfunctiondefinition}
	\end{align}
\end{Def}

\begin{Que} \label[Que]{Que:ackermannfunctiondefinition}
	\(\NaturalNumbers^2\)上に，\(\pair{x_1, y_1} \prec \pair{x_2, y_2}\)であることを以下のいずれかが満たされることと定義して\(\NaturalNumbers^2\)上の二項関係を定める：
	\begin{itemize}
		\item \(x_1 < x_2\)である．
		\item \(x_1 = x_2\)であり，かつ\(y_1 < y_2\)である．
	\end{itemize}
	このとき，\(\prec\)は\(\NaturalNumbers^2\)上の整礎な半順序である．
	Ackermann関数の定義を\(\NaturalNumbers^2\)をこの順序\(\prec\)による整礎帰納法によって正当化せよ．
	すなわち，\cref{eq:ackermannfunctiondefinition}を満たす関数\(\Ack \colon \NaturalNumbers^2 \to \NaturalNumbers\)がただ1つ存在することを示せ．
\end{Que}

\begin{Ex} \label[Ex]{Ex:constantprimitiverecursive}
	\(c\)を自然数とする．このとき，正の整数\(n\)に対して\(\NaturalNumbers^n\)から\(\NaturalNumbers\)への定数関数\(x \mapsto c\)は原始帰納的関数である．
	これを示そう．原始帰納的関数全体の集合を\(\symcal{F}\)として，以下のように写像\(F \colon \NaturalNumbers \to \symcal{F}\)を帰納的に定義する：
	\begin{align*}
		\apply{F}{0}     & = \zero^1,                                                                  \\
		\apply{F}{n + 1} & = \suc \circ \paren{\apply{F}{n}} \quad \text{(\(n \in \NaturalNumbers\))}.
	\end{align*}
	この\(F\)に対し，\(\apply{F}{c}\)は定数関数\(x \mapsto c\)である．
\end{Ex}

\begin{Ex} \label[Ex]{Ex:identitypromitiverecursive}
	\(\NaturalNumbers\)上の恒等写像\(\iota \colon x \mapsto x\)は原始帰納的関数である．
	実際，関数\(\iota_0 \colon \NaturalNumbers^2 \to \NaturalNumbers\)を以下のように原始帰納法によって定義すると，恒等写像\(\iota\)は\(\iota_0\)と
	射影関数\(\proj_2^2\)の合成によって定義できる．
	\begin{align*}
		\apply{\iota_0}{x, 0}     & = \apply{\zero}{x, 0} \quad \text{(\(x \in \NaturalNumbers\))},    \\
		\apply{\iota_0}{x, y + 1} & = \apply{\suc}{{\apply{\proj_2^3}{{x, y, \apply{\iota_0}{x, y}}}}}
		\quad \text{(\(x, y \in \NaturalNumbers\))}.
	\end{align*}
	ここで，規則5における関数\(g, h\)としてそれぞれ\(\zero^2\)と\(\suc \circ \proj_2^3\)をとった．
\end{Ex}

\begin{Ex} \label[Ex]{Ex:addition}
	加法を表す\(\NaturalNumbers^2\)から\(\NaturalNumbers\)への関数\(\pair{x_1, x_2} \mapsto x_1 + x_2\)は原始帰納的関数である．
	実際，以下のように原始帰納法によって定義される関数\(p \colon \NaturalNumbers^2 \to \NaturalNumbers\)が加法を表す関数に相当する：
	\begin{align*}
		\apply{p}{x, 0}     & = \apply{\proj_1^2}{x, 0}
		\quad \text{(\(x \in \NaturalNumbers\))},                                      \\
		\apply{p}{x, y + 1} & = \apply{\suc}{\apply{\proj_3^3}{x, y, \apply{p}{x, y}}}
		\quad \text{(\(x, y \in \NaturalNumbers\))}.
	\end{align*}
\end{Ex}

\begin{Ex} \label[Ex]{Ex:predecessorprimitiverecursive}
	自然数\(x\)の前者\(x - 1\)を表す関数\(m \colon \NaturalNumbers \to \NaturalNumbers\)は原始帰納的関数である．
	実際，\(m\)は原始帰納法によって以下のように定義される関数\(m_0\)と\(\proj_2^2\)との合成によって定義できる：
	\begin{align*}
		\apply{m_0}{x, 0}     & = \apply{\zero^2}{x, 0}	\quad \text{(\(x \in \NaturalNumbers\))}, \\
		\apply{m_0}{x, y + 1} & = \apply{\proj_2^3}{x, y, \apply{m_0}{x}}
		\quad \text{(\(x, y \in \NaturalNumbers\))}.
	\end{align*}

	「\(\apply{m}{x}\)で\(x\)の前者を表す」という意図に沿って考えると，
	本来\(\apply{m}{0}\)は\(0-1\)となるべきだが，これは自然数ではないので便宜的にその値を\(0\)としている．
\end{Ex}

\begin{Ex} \label[Ex]{Ex:minusprimitiverecursivefunction}
	減法を表す\(\NaturalNumbers^2\)から\(\NaturalNumbers\)への以下の関数\(\entireminus\)は原始帰納的関数である．
	実際，以下のように原始帰納法によって定義される関数\(d \colon  \NaturalNumbers^2 \to \NaturalNumbers\)が減法を表す関数に相当する：
	\begin{align*}
		\apply{d}{x, 0}     & = \apply{\proj_1^2}{x}
		\quad \text{(\(x \in \NaturalNumbers\))},                                   \\
		\apply{d}{x, y + 1} & = \apply{\proj_3^3}{x, y, \apply{m}{\apply{d}{x, y}}}
		\quad \text{(\(x, y \in \NaturalNumbers\))}.
	\end{align*}
	ここで，\(m \colon \NaturalNumbers \to \NaturalNumbers\)は\cref{Ex:predecessorprimitiverecursive}における前者を表す関数である．

	\(x, y \in \NaturalNumbers\)に対する\(\apply{d}{x, y}\)を
	\begin{equation}
		x \entireminus y
		\label{eq:entireminus}
	\end{equation}
	と表す．
	\(x \geq y\)ならば\(x \entireminus y = x - y\)であり（つまり通常の減法と同じ），\(x < y\)ならば\(x \entireminus y = 0\)である．
\end{Ex}

\begin{Que} \label[Que]{Que:primitiverevursivefunction}
	\(n\)を正の整数とする．
	以下の\(\NaturalNumbers, \NaturalNumbers^2, \NaturalNumbers^n\)から\(\NaturalNumbers\)への関数はいずれも原始帰納的関数であることを示せ．
	\begin{align*}
		\pair{x, y}                 & \mapsto xy,                              \\
		\pair{x, y}                 & \mapsto x^y,                             \\
		x                           & \mapsto x!,                              \\
		\pair{x_1, x_2, \dots, x_n} & \mapsto \max \Set{x_1, x_2, \dots, x_n}, \\
		\pair{x_1, x_2, \dots, x_n} & \mapsto \min \Set{x_1, x_2, \dots, x_n}.
	\end{align*}
	ただし，\(0! = 0^0 = 1\)と決めておくこととする．
\end{Que}

\index[sidx]{\(\truefunction\)：真であることを表す関数}
\index[sidx]{\(\falsefunction\)：偽であることを表す関数}
\begin{Que} \label[Que]{Que:truefalseprimitiverecursive}
	以下で定義される関数\(\truefunction, \falsefunction \colon \NaturalNumbers \to \NaturalNumbers\)が原始帰納的関数であることを示せ：
	\begin{align}
		\apply{\truefunction}{x}  & =
		\begin{cases*}
			0 & (\(x = 0\)のとき), \\
			1 & (それ以外のとき),
		\end{cases*}
		\label{eq:trueprimitiverecursive} \\
		\apply{\falsefunction}{x} & =
		\begin{cases*}
			1 & (\(x = 0\)のとき), \\
			0 & (それ以外のとき).
		\end{cases*}
		\label{eq:falseprimitiverecursive}
	\end{align}
\end{Que}

上記の例や\cref{Que:primitiverevursivefunction}により，以下の主張が従う．

\begin{Lemma} \label[Lemma]{lemma:recursivefunctionoperation}
	原始帰納的関数の和，（\(\entireminus\)の意味での）差，積，べき乗によって表される関数はすべて原始帰納的関数であり，
	帰納的関数の和，（\(\entireminus\)の意味での）差，積，べき乗によって表される関数はすべて帰納的関数である．
\end{Lemma}

\begin{Que} \label[Que]{Que:absolutedifferenceprimitiverecursive}
	以下の関数は原始帰納的関数であることを示せ．
	\[
		\pair{x, y} \mapsto \absolute{x - y}
	\]
\end{Que}

\begin{Que} \label[Que]{Que:remquoprimitiverecursive}
	以下の関数は原始帰納的関数であることを示せ．
	\begin{enumerate}
		\item \(x, y \in \NaturalNumbers\)に対し，\(y = 0\)ならば\(x\)を，\(y \neq 0\)ならば\(x\)を\(y\)で割った余りをそれぞれ対応させる関数\(\rem \colon \NaturalNumbers^2 \to \NaturalNumbers\)，
		\item \(x, y \in \NaturalNumbers\)に対し，\(y = 0\)ならば0を\(y \neq 0\)ならば\(x\)を\(y\)で割った商をそれぞれ対応させる\(\quo \colon \NaturalNumbers^2 \to \NaturalNumbers\)．
	\end{enumerate}
\end{Que}

以上の事実から，\cref{sec:codizeforfinitesequence}で定義したCantorの対関数や有限列のコード化に関する種々の関数が原始帰納的関数であることがわかる．

\begin{Lemma} \label[Lemma]{lemma:codizeforfinitesequenceprimitiverecursive}
	\Cref{sec:codizeforfinitesequence}で定義した関数\(\pairfunc, \pairfunc^n, \pairfunc_i^n, \connectionfunc\)はすべて原始帰納的関数である．
	また，固定された自然数\(n\)に対して，
	長さ\(n\)の自然数からなる有限列\(x_1, x_2, \dots, x_n\)をそのコード
	\(\finitesequencecode{x_1, x_2, \dots, x_n}\)に対応させる\(\NaturalNumbers^n\)から\(\NaturalNumbers\)への関数も原始帰納的関数である．
	さらに，自然数\(i, x\)を\(\codefunctioninv{x}{i}\)に対応させる\(\NaturalNumbers^2\)から\(\NaturalNumbers\)への関数と
	自然数\(x, y\)を\(x \codeconnection y\)に対応させる\(\NaturalNumbers^2\)から\(\NaturalNumbers\)への関数も原始帰納的関数である．
\end{Lemma}

\begin{Def} \label{Def:pisumprimitiverecursive}
	\(n\)を正の整数とする．
	関数\(\varphi \colon \NaturalNumbers^{n+1} \to \NaturalNumbers\)が与えられたとき，
	\(x_1, x_2, \dots, x_n, z \in \NaturalNumbers\)に対して
	\begin{align*}
		\apply{\varphi}{x_1, x_2, \dots, x_n, 0}
		+ \apply{\varphi}{x_1, x_2, \dots, x_n, 1}
		+ \dots
		+ \apply{\varphi}{x_1, x_2, \dots, x_n, z \entireminus 1}, \\
		\apply{\varphi}{x_1, x_2, \dots, x_n, 0}
		\times \apply{\varphi}{x_1, x_2, \dots, x_n, 1}
		\times \dots
		\times \apply{\varphi}{x_1, x_2, \dots, x_n, z \entireminus 1}
	\end{align*}
	をそれぞれ
	\begin{align}
		\sum_{y < z} \apply{\varphi}{x_1, x_2, \dots, x_n, y},
		\label{eq:sumprimitiverecursive} \\
		\prod_{y < z} \apply{\varphi}{x_1, x_2, \dots, x_n, y}
		\label{eq:piprimitiverecursive}
	\end{align}
	と表す．ただし，\(z = 0\)のときは\cref{eq:sumprimitiverecursive}と\cref{eq:piprimitiverecursive}の値はそれぞれ
	\(0, 1\)であるものとする．
\end{Def}

\Cref{lemma:recursivefunctionoperation}により，以下の定理が得られる．
\begin{Lemma} \label[Lemma]{lemma:sumprimitiverecursive}
	\(\varphi \colon \NaturalNumbers^{n+1} \to \NaturalNumbers\)が原始帰納的関数であるならば
	\(\pair{x_1, x_2, \dots, x_n, z} \in \NaturalNumbers^{n + 1}\)を\cref{eq:sumprimitiverecursive}, \cref{eq:piprimitiverecursive}に対応させる
	\(\NaturalNumbers^{n+1}\)から\(\NaturalNumbers\)への関数はいずれも原始帰納的関数であり，
	\(\varphi\)が帰納的関数であるならば
	\(\pair{x_1, x_2, \dots, x_n, z} \in \NaturalNumbers^{n + 1}\)を\cref{eq:sumprimitiverecursive}, \cref{eq:piprimitiverecursive}に対応させる
	\(\NaturalNumbers^{n+1}\)から\(\NaturalNumbers\)への関数はいずれも帰納的関数である．
\end{Lemma}

\index[widx]{げんしきのうてきじゅつご@原始帰納的述語}
\index[widx]{きのうてきじゅつご@帰納的述語}
\begin{Def} \label{Def:recursivepredicate}
	\(n\)を正の整数とする．\(\NaturalNumbers\)上の\(n\)項関係\(R \subset \NaturalNumbers^n\)が
	\term{原始帰納的述語}であるとは，
	以下で定義される\(R\)の特性関数\(\characteristicfunction{R} \colon \NaturalNumbers^n \to \NaturalNumbers\)が原始帰納的関数であることをいう：
	\begin{align}
		\apply{\characteristicfunction{R}}{x_1, x_2, \dots, x_n} =
		\begin{cases*}
			1 & (\(\pair{x_1, x_2, \dots, x_n} \in R\)のとき), \\
			0 & (それ以外のとき).
		\end{cases*}
		\label{eq:characteristicfunction}
	\end{align}
	また，\(\characteristicfunction{R}\)が帰納的関数であるならば，\(R\)は\term{帰納的述語}であるという．
\end{Def}

\begin{Ex} \label[Ex]{Ex:recursivepredicate}
	二項関係\(x = y, x < y\)は原始帰納的述語である．実際，それぞれの特性関数が
	\begin{align*}
		\characteristicfunction{x = y} & = \apply{\falsefunction}{\absolute{x-y}}  \\
		\characteristicfunction{x < y} & = \apply{\truefunction}{y \entireminus x}
	\end{align*}
	と表せる．
\end{Ex}

\begin{Lemma} \label[Lemma]{lemma:relationoperationrecursive}
	\(n\)を正の整数とする．\(R \subset \NaturalNumbers^n\)が原始帰納的述語であるならば，
	\[
		\complementset{R}
	\]
	も原始帰納的述語である．ここで，\(\complementset{R} = \NaturalNumbers^n \setminus R\)は\(R\)の補集合である．
	また，\(R_1, R_2 \subset \NaturalNumbers^n\)が原始帰納的述語であるならば，
	\begin{align*}
		 & R_1 \cup R_2, \\
		 & R_1 \cap R_2
	\end{align*}
	も原始帰納的述語である．
	上記「原始帰納的」を「帰納的」に置き換えた主張も同様に成り立つ．
\end{Lemma}

\begin{proof}
	それぞれの特性関数が以下のように書けることから従う：
	\begin{align*}
		\apply{\characteristicfunction{\complementset{R}}}{x} & = \apply{\falsefunction}{\apply{\characteristicfunction{R}}{x}}
		\quad \text{(\(x \in \NaturalNumbers^n\))},                                                                                                                        \\
		\apply{\characteristicfunction{R_1 \cup R_2}}{x}      & = \apply{\truefunction}{\apply{\characteristicfunction{R_1}}{x} + \apply{\characteristicfunction{R_2}}{x}}
		\quad \text{(\(x \in \NaturalNumbers^n\))},                                                                                                                        \\
		\apply{\characteristicfunction{R_1 \cap R_2}}{x}      & = \apply{\characteristicfunction{R_1}}{x} \cdot \apply{\characteristicfunction{R_2}}{x}
		\quad \text{(\(x \in \NaturalNumbers^n\))}.
	\end{align*}
\end{proof}

\begin{Lemma} \label[Lemma]{lemma:casesrecursive}
	\(n, m\)を正の整数とする．
	\(R_1, R_2, \dots, R_m \subset \NaturalNumbers^n\)
	を，どの2つも同時に成り立たない，つまり\(i \neq j\)ならば\(R_i \cap R_j = \emptyset\)であるような\(\NaturalNumbers\)上の\(n\)項関係とする．
	関数\(\varphi_1, \varphi_2, \dots, \varphi_m, \varphi_{m+1} \colon \NaturalNumbers^n \to \NaturalNumbers\)に対して，
	関数\(\varphi \colon \NaturalNumbers^n \to \NaturalNumbers\)を
	\begin{align*}
		\apply{\varphi}{x_1, x_2, \dots, x_n} =
		\begin{cases*}
			\apply{\varphi_1}{x_1, x_2, \dots, x_n}     & (\(\pair{x_1, \dots, x_n} \in R_1\)のとき), \\
			\apply{\varphi_2}{x_1, x_2, \dots, x_n}     & (\(\pair{x_1, \dots, x_n} \in R_2\)のとき), \\
			\dotsb                                      &                                          \\
			\apply{\varphi_m}{x_1, x_2, \dots, x_n}     & (\(\pair{x_1, \dots, x_n} \in R_m\)のとき), \\
			\apply{\varphi_{m+1}}{x_1, x_2, \dots, x_n} & (上記以外のとき)
		\end{cases*}
	\end{align*}
	と定義する．このとき，以下が成り立つ．
	\begin{enumerate}
		\item 各\(R_i\)がすべて原始帰納的述語で，各\(\varphi_i\)がすべて原始帰納的関数であるならば，\(\varphi\)も原始帰納的関数である．
		\item 各\(R_i\)がすべて帰納的述語で，各\(\varphi_i\)がすべて帰納的関数であるならば，\(\varphi\)も帰納的関数である．
	\end{enumerate}
\end{Lemma}

\begin{proof}
	\(\varphi\)が
	\begin{align*}
		\apply{\varphi}{x_1, x_2, \dots, x_n} = \sum_{i < m+1} \apply{\characteristicfunction{R_i}}{x_1, x_2, \dots, x_n} \cdot \apply{\varphi_i}{x_1, x_2, \dots, x_n} \\
		\quad + \apply{\characteristicfunction{\complementset{\paren{R_1 \cup R_2 \cup \dotsb \cup R_m}}}}{x_1, x_2, \dots, x_n} \cdot \apply{\varphi_{m+1}}{x_1, x_2, \dots, x_n}
	\end{align*}
	と書けることから従う．
\end{proof}

\begin{Lemma} \label[Lemma]{lemma:forallexistsrecursive}
	\(n\)を自然数とする．\(R\)を\(\NaturalNumbers\)上の\(n+1\)項関係とするとき，
	以下で定義される\(n\)項関係\(R_1, R_2\)を考える．
	\begin{align*}
		                & \pair{x_1, x_2, \dots, x_n, z} \in R_1                                       \\
		\metaequivalent & \text{\(\pair{x_1, x_2, \dots, x_n, y} \in R\)かつ\(y < z\)となる\(y\)が存在する},     \\
		                & \pair{x_1, x_2, \dots, x_n, z} \in R_1                                       \\
		\metaequivalent & \text{\(y < z\)となるすべての\(y\)に対して\(\pair{x_1, x_2, \dots, x_n, y} \in R\)となる}.
	\end{align*}
	このとき，\(R\)が原始帰納的述語であるならば\(R_1, R_2\)も原始帰納的述語であり，
	\(R\)が帰納的述語ならば\(R_1, R_2\)も帰納的述語である．
\end{Lemma}

\begin{proof}
	特性関数がそれぞれ
	\begin{align*}
		 & \apply{\truefunction}{\sum_{y < z} \apply{\characteristicfunction{R}}{x_1, x_2, \dots, x_n, y}}, \\
		 & \apply{\truefunction}{\prod_{y < z} \apply{\characteristicfunction{R}}{x_1, x_2, \dots, x_n, y}} \\
	\end{align*}
	と表されることから従う．
\end{proof}

\index[sidx]{\(\muoperator_{y<z}\)：有界\(\muoperator\)演算子}
\index[widx]{\(\muoperator\)さようそ@\(\muoperator\)さようそ!ゆうかい\(\muoperator\)作用素@有界---}
\begin{Def} \label{Def:boundedmuoperator}
	\(n\)を正の整数とし，\(R\)を\(\NaturalNumbers\)上の\(n\)項関係とする．
	各\(x_1, x_2, \dots, x_n, z \in \NaturalNumbers\)に対して，
	\(y < z\)かつ\(\pair{x_1, x_2, \dots, x_n, y} \in R\)となる\(y \in \NaturalNumbers\)が存在するならば
	そのような\(y\)の最小値を，存在しないならば\(z\)を
	\begin{equation}
		\muoperator_{y < z} \apply{R}{x_1, x_2, \dots, x_n, y}
		\label{eq:boundedmuoperator}
	\end{equation}
	と表す．
\end{Def}

\begin{Lemma} \label[Lemma]{lemma:boundedmuoperatorrecursive}
	\(R\)が原始帰納的述語であれば
	\(\pair{x_1, x_2, \dots, x_n, z}\)に対して\cref{eq:boundedmuoperator}で定義される自然数を対応させる関数は原始帰納的関数である．
	\(R\)が帰納的述語であれば
	\(\pair{x_1, x_2, \dots, x_n, z}\)に対して\cref{eq:boundedmuoperator}で定義される自然数を対応させる関数は帰納的関数である．
\end{Lemma}

\begin{proof}
	\[
		\muoperator_{y<z} \apply{R}{x_1, x_2, \dots, x_n, y} = \sum_{w < z} \prod_{y < w} \apply{\characteristicfunction{\complementset{R}}}{x_1, x_2, \dots, x_n, y}
	\]
	と書けることから従う．
\end{proof}


以上の議論では，帰納的関数や原始帰納的関数は\(\NaturalNumbers^n\)上全域で定義された関数を対象にした．
ここでは，必ずしも\(\NaturalNumbers^n\)上全域で定義されているとは限らない関数について考えよう．
まずは以下のように記号を定義しておく．

\index[sidx]{\(\partialequal\)：部分関数に対する等号の拡張}
\index[widx]{ぶぶんかんすう@部分関数}
\begin{Def} \label{Def:partialequal}
	\(n\)を正の整数とする．\(X, Y \subset \NaturalNumbers^n\)に対して，
	関数\(f \colon X \to \NaturalNumbers\)と関数\(g \colon Y \to \NaturalNumbers\)を考える．
	このような\(f, g\)を\(\NaturalNumbers\)の\term{部分関数}という．
	このとき，各\(x_1, x_2, \dots, x_n \in \NaturalNumbers\)に対して
	\begin{equation}
		\apply{f}{x_1, x_2, \dots, x_n} \partialequal \apply{g}{x_1, x_2, \dots, x_n}
		\label{eq:partialequal}
	\end{equation}
	を，以下の表に従って定める．

	\begin{table}[htbp]
		\centering
		\caption{\(\apply{f}{x_1, x_2, \dots, x_n} \partialequal \apply{g}{x_1, x_2, \dots, x_n}\)の真偽}
		\label{tab:partialequal}
		\begin{tabular}{cccc}
			\toprule
			\(\pair{x_1, x_2, \dots, x_n} \in X\) &
			\(\pair{x_1, x_2, \dots, x_n} \in Y\) &
			真偽                                                        \\
			\midrule
			成り立つ                                  & 成り立つ   & 通常の等号と同じ \\
			成り立たない                                & 成り立つ   & 偽        \\
			成り立つ                                  & 成り立たない & 偽        \\
			成り立たない                                & 成り立たない & 真        \\
			\bottomrule
		\end{tabular}
	\end{table}

	ここで，\cref{tab:partialequal}における「通常の等号と同じ」とは，その真偽が
	\[
		\apply{f}{x_1, x_2, \dots, x_n} = \apply{g}{x_1, x_2, \dots, x_n}
	\]
	が成り立つことと同じであるという意味である．

	また，\(X \subset \NaturalNumbers^n\)とするとき，\(X\)の部分集合\(R \subset X\)を\(n\)個の自然数の間の関係を表したものとみなすとき，
	\(R\)を\(\NaturalNumbers^n\)上の\term{部分述語}という．
	このとき，\(R\)の特性関数\(\characteristicfunction{R} \colon X \to \NaturalNumbers\)は，
	\cref{eq:characteristicfunction}における\(\NaturalNumbers^n\)を\(X\)に置き換えることによって得られる．
\end{Def}

さて，上記の概念をもとにして定義域が\(\NaturalNumbers^n\)全域とは限らない場合の帰納的関数と原始帰納的関数を定義しよう．

\index[widx]{きのうてきぶぶんかんすう@帰納的部分関数}
\index[widx]{げんしきのうてきぶぶんかんすう@原始帰納的部分関数}
\begin{Def} \label{Def:partialrecursivefunction}
	\term{帰納的部分関数}は，以下のように帰納的に定義される：
	\begin{enumerate}
		\item ゼロ関数，後者関数，射影関数はいずれも帰納的部分関数である．
		\item 合成と原始帰納法については，\cref{Def:recursivefunction}における等号を\(\partialequal\)に置き換えたもので同様に定義する．
		\item \(g\)が\(\NaturalNumbers^{n+1}\)上の帰納的部分関数であるとき，
		      \begin{equation}
			      \apply{f}{x_1, x_2, \dots, x_n} \partialequal \muoperator y \paren{\apply{g}{x_1, x_2, \dots, x_n, y} = 0}
			      \label{eq:partialequalminimization}
		      \end{equation}
		      で定義される\(\NaturalNumbers^n\)上の部分関数\(f\)は帰納的部分関数である．
		      ここで，\(\muoperator y \paren{\apply{g}{x_1, x_2, \dots, x_n, y} = 0}\)の値が定義されており，
		      その値が\(a\)であるとは，すべての\(b < a\)に対して\(\apply{g}{x_1, x_2, \dots, x_n, b}\)の値が定義されており，
		      その値が0より大きく，しかも\(\apply{g}{x_1, x_2, \dots, x_n, a} = 0\)であることをいう．
	\end{enumerate}

	また，上記規則3を一度も適用せずに得られる帰納的部分関数をとくに
	\term{原始帰納的部分関数}という．

	さらに，\(\NaturalNumbers^n\)上の部分述語\(R\)が\term{帰納的部分述語}であるとは，
	その特性関数が帰納的部分関数であることをいい，\(R\)が\term{原始帰納的部分述語}であるとは，
	その特性関数が帰納的部分関数であることをいう．
\end{Def}

\Cref{Def:partialrecursivefunction}では，最小化において最小値の存在を仮定しなかった．
これは，\(\apply{f}{x_1, x_2, \dots, x_n}\)の値が定義されている必要がなくなったことによる．

帰納的関数において重要な結果となるのが以下に述べるKleeneの標準形定理である．
証明には長い準備が必要なため，廣瀬\cite{hirose2024}を参照のこと．

\index[widx]{Kleeneのひょうじゅんけいていり@Kleeneの標準形定理}
\begin{Thm}[Kleeneの標準形定理]
	\(n\)を正の整数とする．原始帰納的関数\(U \colon \NaturalNumbers \to \NaturalNumbers\)と
	\(\NaturalNumbers^{n+1}\)が存在して，
	任意の\(\NaturalNumbers^n\)上の帰納的部分関数は，自然数\(e\)を用いて
	\begin{equation}
		\apply{f}{x_1, x_2, \dots, x_n} \partialequal \apply{U}{\muoperator y \apply{T_n}{e, x_1, x_2, \dots, x_n, y}}
		\quad \text{(\(x_1, x_2, \dots, x_n \in \NaturalNumbers\))}
		\label{eq:kleenenormaltheorem}
	\end{equation}
	と表せる．ここで，\(\muoperator y \apply{T_n}{e, x_1, x_2, \dots, x_n, y}\)は，
	\(\pair{e, x_1, x_2, \dots, x_n, y} \in T_n\)が成り立つような\(y\)のうち最小のものである．
\end{Thm}

\section{計算可能関数} \label[section]{sec:computablefunction}

帰納的関数は，今でこそ計算可能性の定式化の1つであるとみなされているが，
定義だけを見ると「計算」というニュアンスはあまり感じられないかもしれない．
そこで，もう1つの定式化として形式的体系によるものを紹介する．
廣瀬\cite{hirose2024}に従い，この体系を\(\symcal{R}\)と表記する．
形式的体系というと\cref{chap:proof}で紹介した自然演繹式シークエント計算を思い浮かべるかもしれないが，
それとはまったく別の体系である．
自然演繹式シークエント計算は数学における「論理」をシミュレートするために導入した体系であったが，
今回導入する形式的体系は，数学における「計算」をシミュレートするための体系である．

体系の内容としては計算機科学における操作的意味論によるプログラミング言語の動作定義に近い．
例えば，古典ではあるがいわゆるTaPL\cite{pierce2013types_jp}を見れば，おおよそ似たような形で式の評価規則が定義されている．

まったく別の体系とは言っても，使う記号や定義は\cref{chap:syntax}で紹介したものと非常に似通っている．
ただし，異なる点も多いので，イチからすべて説明しよう．

まずは使う記号である．\Cref{tab:computablefunctionsymbol}に体系\(\symcal{R}\)で使用する記号の一覧を示す．

\index[widx]{へんすうきごう@変数記号}
\index[widx]{かんすうきごう@関数記号}
\index[widx]{ていすうきごう@定数記号}
\index[widx]{かんけいきごう@関係記号}
\index[widx]{ほじょきごう@補助記号}
\index[sidx]{\(\symcal{R}\)：計算可能関数を定義する形式的体系}
\index[sidx]{\(\obj{0}\)：体系\(\symcal{R}\)における定数記号}
\index[sidx]{\(\obj{\succfunc}\)：体系\(\symcal{R}\)における関数記号}
\begin{table}[htbp]
	\centering
	\caption{体系\(\symcal{R}\)で使用する記号の一覧}
	\label{tab:computablefunctionsymbol}
	\begin{tabular}{ccc}
		\toprule
		種別        & 一覧                                                                   & 備考              \\
		\midrule
		変数記号      & \(\obj{x}, \obj{y}, \obj{z}, \dotsc\)                                & 無限に多く存在する（可算無限） \\
		関数記号      & \(\obj{f}, \obj{g}, \obj{h}, \dotsc\)                                & 無限に多く存在する（可算無限） \\
		（特定の）定数記号 & \(\obj{0}\)                                                          &                 \\
		（特定の）関数記号 & \(\obj{\succfunc}\)                                                  &                 \\
		補助記号      & \(\text{``\(\lparen\)''}, \text{``\(\rparen\)''}, \text{``\(,\)''}\) & カッコやカンマ         \\
		\bottomrule
	\end{tabular}
\end{table}

1階述語論理のときと似ているが少し違う．例えば1階述語論理では言語という形で定数記号や関数記号にバリエーションを持たせることができたが，
\(\symcal{R}\)では定数記号は\(\obj{0}\)のみ，関数記号は特定の\(\obj{\succfunc}\)の他は変数記号と似たような取り扱いがなされている．
これは，\(\symcal{R}\)は計算可能関数を定義するための体系であることからくる違いである．

さて，次は項である．定義は\cref{Def:term}とほぼ同じである．

\begin{Def} \label{Def:computablefunctionterm}
	\(\symcal{R}\)における項は，以下のように帰納的に定義される．
	\begin{enumerate}
		\item 変数記号はすべて項である．
		\item 定数記号\(\obj{0}\)は項である．
		\item 項\(t_1, t_2, \dotsc, t_n\)と関数記号\(f\)に対して，\(\apply{f}{t_1, t_2, \dots, t_n}\)も項である．
		\item 以上の規則を有限回適用して得られる記号列のみが項である．
	\end{enumerate}
\end{Def}

1階述語論理では関数記号は言語側が有していたためアリティの存在を要請していたが，
\(\symcal{R}\)ではアリティの存在を要請しない．
関数記号を無限に多く存在することを要請しているのは「すでに登場した関数記号のいずれとも異なる関数記号をいつでも用意できる」ことを担保するためである．
ここで関数記号に対してアリティの存在を要請してしまうと「今欲しいアリティの関数記号がもうない」という事態に陥りかねない．
そこで，\(\symcal{R}\)では関数記号に対してアリティの存在を要請せず，項を組み立てる際に組み立てる側が自由にアリティを指定できるようにしている．

項の略記法を定義しておく．

\begin{Def} \label{Def:computablefunctionnumeral}
	\(n\)を自然数とする．項\(\numeral{n}\)を以下のように帰納的に定義する．
	\begin{align}
		\numeral{0}     & = \obj{0},                                                                       \\
		\numeral{n + 1} & = \apply{\obj{\succfunc}}{\numeral{n}} \quad \text{(\(n \in \NaturalNumbers\))}.
	\end{align}
	自然数\(n\)に対する項\(\numeral{n}\)を，\(n\)に対する\term{数項}と呼ぶ．
\end{Def}

計算結果を示すための記号列を定義しよう．
\Cref{Def:evaluation}は，\cref{chap:proof}におけるシークエントの定義（\cref{Def:sequent}）に相当する．

\index[widx]{ひょうかしき@評価式}
\index[sidx]{\(t_1 \evaluate t_2\)：評価式}
\index[widx]{ひょうかしき@評価式!うへん@---の右辺}
\index[widx]{ひょうかしき@評価式!さへん@---の左辺}
\begin{Def} \label{Def:evaluation}
	項\(t_1, t_2\)を用いて
	\begin{equation}
		t_1 \evaluate t_2
		\label{eq:evaluation}
	\end{equation}
	と表される記号列を\term{評価式}と呼ぶ．
	\(t_1\)をこの評価式の\term{左辺}，\(t_2\)をこの評価式の\term{右辺}と呼ぶ．
\end{Def}

次に導入すべきは推論規則である．まずはそのために必要な概念を定義しよう．

\begin{Def} \label{Def:computablefunctionsubstitution}
	\(e\)を評価式，\(a\)を変数記号，\(n\)を数項とする．
	このとき，\(e\)に現れるすべての\(a\)を\(n\)に置き換えた評価式を
	\begin{equation}
		\SubEvaluate{e}{a}{n}
		\label{eq:computablefunctionsubstitution}
	\end{equation}
	と表す．
\end{Def}

\begin{Def} \label{Def:computablefunctionreplacement}
	\(t\)を項，\(f\)を関数記号，\(n_1, n_2, \dots, n_r, n\)を数項とする．
	このとき，\(t\)に現れるすべての\(\apply{f}{n_1, n_2, \dots, n_r}\)を\(n\)に置き換えた項を
	\begin{equation}
		\RepTerm{t}{\apply{f}{n_1, n_2, \dots, n_r}}{n}
		\label{eq:computablefunctionreplacement}
	\end{equation}
	と表す．
\end{Def}

いよいよ推論規則を定義しよう．

\begin{Def} \label{Def:computablefunctioninference}
	\(e\)を評価式，\(a\)を変数記号，\(n\)を数項とする．
	このとき，下記は\(\symcal{R}\)における推論規則である．
	\begin{prooftree}
		\AxiomC{\(e\)}
		\LeftLabel{(SUB)}
		\UnaryInfC{\(\SubEvaluate{e}{a}{n}\)}
	\end{prooftree}

	また，\(t_1, t_2\)を項，\(f\)を関数記号，\(n_1, n_2, \dots, n_r, n\)を数項とするとき，
	下記は\(\symcal{R}\)における推論規則である．
	\begin{prooftree}
		\AxiomC{\(t_1 \evaluate t_2\)}
		\AxiomC{\(\apply{f}{n_1, n_2, \dots, n_r} \evaluate n\)}
		\LeftLabel{(REPL)}
		\BinaryInfC{\(t_1 \evaluate \RepTerm{t_2}{\apply{f}{n_1, n_2, \dots, n_r}}{n}\)}
	\end{prooftree}
\end{Def}

推論規則はこの2つだけである．
ここから評価式の導出が定義できる．
以下の定義は，\cref{chap:proof}における証明図の定義（\cref{Def:proofDiagram}）に相当する．

\index[widx]{どうしゅつず@導出図}
\begin{Def} \label{Def:evaluatable}
	\(E\)を評価式の有限集合とする．評価式\(e\)を下段とする\(E\)からの\term{導出図}を，以下のように帰納的に定義する：
	\begin{enumerate}
		\item \(e \in E\)ならば，以下の図式は\(e\)を下段とする\(E\)からの導出図である：
		      \begin{prooftree}
			      \AxiomC{}
			      \UnaryInfC{\(e\)}
		      \end{prooftree}
		\item \(\symcal{D}\)が評価式\(e\)を下段とする\(E\)からの導出図であり，\(a\)が変数記号，\(n\)が数項であるならば，
		      以下の図式は\(\SubEvaluate{e}{a}{n}\)を下段とする\(E\)からの導出図である：
		      \begin{prooftree}
			      \AxiomC{\(\symcal{D}\)}
			      \LeftLabel{(SUB)}
			      \UnaryInfC{\(\SubEvaluate{e}{a}{n}\)}
		      \end{prooftree}
		\item \(\symcal{D}_1\)が評価式\(t_1 \evaluate t_2\)を下段とする\(E\)からの導出図であり，
		      \(\symcal{D}_2\)が評価式\(\apply{f}{n_1, n_2, \dots, n_r} \evaluate n\)を下段とする\(E\)からの導出図であるならば，
		      以下の図式は評価式\(t_1 \evaluate \RepTerm{t_2}{\apply{f}{n_1, n_2, \dots, n_r}}{n}\)を下段とする\(E\)からの導出図である：
		      \begin{prooftree}
			      \AxiomC{\(\symcal{D}_1\)}
			      \AxiomC{\(\symcal{D}_2\)}
			      \LeftLabel{(REPL)}
			      \BinaryInfC{\(t_1 \evaluate \RepTerm{t_2}{\apply{f}{n_1, n_2, \dots, n_r}}{n}\)}
		      \end{prooftree}
		\item 以上の規則を有限回適用して得られるもののみが\(E\)からの導出図である．
	\end{enumerate}
\end{Def}

導出図とくればあとは導出可能性である．
\Cref{chap:proof}における証明可能性の定義（\cref{Def:provable}）を思い起こせば，
以下の定義は極めて自然であろう．

\index[widx]{どうしゅつかのう@導出可能}
\index[sidx]{\(E \evaluatable e\)：\(E\)から導出可能}
\begin{Def} \label{Def:computablefunctionevaluatable}
	\(E\)を評価式の有限集合とする．評価式\(e\)が\(E\)から\term{導出可能}であるとは，
	\(e\)を下段とする\(E\)からの導出図が存在することをいう．
	このことを
	\begin{equation}
		E \evaluatable e
		\label{eq:computablefunctionevaluatable}
	\end{equation}
	と表す．
\end{Def}

\begin{Ex} \label[Ex]{Ex:computablefunctionevaluatable}
	\(E\)を以下の2つの評価式からなる集合とする：
	\begin{align*}
		\apply{\obj{f}}{a_0, \numeral{0}}            & \evaluate a_0,                                          \\
		\apply{\obj{f}}{a_0, \apply{\succfunc}{a_1}} & \evaluate \apply{\succfunc}{\apply{\obj{f}}{a_0, a_1}}.
	\end{align*}
\(x\)を自然数とする．このとき，以下は評価式\(\apply{\obj{f}}{\numeral{x}, \numeral{1}} \evaluate \apply{\succfunc}{\numeral{x}}\)の\(E\)からの導出図である：
	\begin{prooftree}
		\AxiomC{}
		\UnaryInfC{\(\apply{\obj{f}}{a_0, \numeral{0}} \evaluate a_0\)}
		\LeftLabel{(SUB)}
		\UnaryInfC{\(\apply{\obj{f}}{\numeral{x}, \numeral{0}} \evaluate \numeral{x}\)}
		\AxiomC{}
		\UnaryInfC{\(\apply{\obj{f}}{a_0, \apply{\succfunc}{a_1}} \evaluate \apply{\succfunc}{\apply{\obj{f}}{a_0, a_1}}\)}
		\LeftLabel{(SUB)}
		\UnaryInfC{\(\apply{\obj{f}}{a_0, \numeral{1}} \evaluate \apply{\succfunc}{\apply{\obj{f}}{a_0, \numeral{0}}}\)}
		\LeftLabel{(SUB)}
		\UnaryInfC{\(\apply{\obj{f}}{\numeral{x}, \numeral{1}} \evaluate \apply{\succfunc}{\apply{\obj{f}}{\numeral{x}, \numeral{0}}}\)}
		\LeftLabel{(REPL)}
		\BinaryInfC{\(\apply{\obj{f}}{\numeral{x}, \numeral{1}} \evaluate \apply{\succfunc}{\numeral{x}}\)}
	\end{prooftree}
	(REPL)の第2の前提として用いた\(\apply{\obj{f}}{\numeral{x}, \numeral{0}} \evaluate \numeral{x}\)は，
	関数記号\(\obj{f}\)の引数がどちらも数項（\(\numeral{x}, \numeral{0}\)）になってはじめて使えることに注意しよう．
	そのため，2つ目の枝でも\(a_1\)だけでなく\(a_0\)にも(SUB)を適用し，\(\apply{\obj{f}}{a_0, \numeral{1}} \evaluate \dots\)の段階でいったん止めずに\(\apply{\obj{f}}{\numeral{x}, \numeral{1}} \evaluate \dots\)まで具体化してから(REPL)を適用している．
\end{Ex}

\Cref{Ex:computablefunctionevaluatable}は，通常の数学において
\begin{align*}
	\apply{f}{x, 0}     & = x \quad (x \in \NaturalNumbers),                     \\
	\apply{f}{x, y + 1} & = \apply{f}{x, y} + 1 \quad (x, y \in \NaturalNumbers)
\end{align*}
と帰納的に定義された関数\(f\)における\(\apply{f}{x, 1}\)の計算を表している．

もう少し複雑な例として，\(E\)からの評価式\(\apply{\obj{f}}{\numeral{2}, \numeral{3}} \evaluate \numeral{5}\)の導出を見てみよう．
\Cref{Thm:lawofexcludedmiddle}の証明と同じように，評価式を書き並べる形で書いてみることとする．

\begin{Ex} \label[Ex]{Ex:computablefunctionevaluatable2}
	\(E\)を\cref{Ex:computablefunctionevaluatable}と同じ集合とする．このとき，
	\(\apply{\obj{f}}{\numeral{2}, \numeral{3}} \evaluate \numeral{5}\)の\(E\)からの導出は以下のように書ける：
	\begin{enumerate}
		\item \(\apply{\obj{f}}{a_0, \numeral{0}} \evaluate a_0\)\quad（\(E\)より）
		\item \(\apply{\obj{f}}{a_0, \apply{\succfunc}{a_1}} \evaluate \apply{\succfunc}{\apply{\obj{f}}{a_0, a_1}}\)\quad（\(E\)より）
		\item \(\apply{\obj{f}}{\numeral{2}, \numeral{0}} \evaluate \numeral{2}\)\quad（1から(SUB)による）
		\item \(\apply{\obj{f}}{\numeral{2}, \numeral{1}} \evaluate \apply{\succfunc}{\apply{\obj{f}}{\numeral{2}, \numeral{0}}}\)\quad（2に(SUB)を2回適用）
		\item \(\apply{\obj{f}}{\numeral{2}, \numeral{1}} \evaluate \numeral{3}\)\quad（3, 4から(REPL)による）
		\item \(\apply{\obj{f}}{\numeral{2}, \numeral{2}} \evaluate \apply{\succfunc}{\apply{\obj{f}}{\numeral{2}, \numeral{1}}}\)\quad（2に(SUB)を2回適用）
		\item \(\apply{\obj{f}}{\numeral{2}, \numeral{2}} \evaluate \numeral{4}\)\quad（5, 6から(REPL)による）
		\item \(\apply{\obj{f}}{\numeral{2}, \numeral{3}} \evaluate \apply{\succfunc}{\apply{\obj{f}}{\numeral{2}, \numeral{2}}}\)\quad（2に(SUB)を2回適用）
		\item \(\apply{\obj{f}}{\numeral{2}, \numeral{3}} \evaluate \numeral{5}\)\quad（7, 8から(REPL)による）
	\end{enumerate}
	各(REPL)の第2の前提（3, 5, 7行目）は，いずれも\(\obj{f}\)の引数がどちらも数項になっている評価式であることに注意しよう．
	\cref{Ex:computablefunctionevaluatable}のときと同様に，(SUB)は\(a_1\)だけでなく\(a_0\)についても毎回適用しており，
	\(a_0\)が\(\numeral{2}\)に定まる前の（変数\(a_0\)を含んだままの）評価式どうしを(REPL)で組み合わせることはしていない．
\end{Ex}

\Cref{Ex:computablefunctionevaluatable2}における評価式\(\apply{\obj{f}}{\numeral{2}, \numeral{3}} \evaluate \numeral{5}\)の導出は，
通常の数学における\(\apply{f}{2, 3}\)の計算に対応している．
使用する規則自体は極めて原始的な代入操作と置換操作のみであるため，まさしく数学における「計算」をシミュレートしているといえるのではないだろうか%
\footnote{%
	\(\symcal{R}\)における導出が数学における「計算」であることは先の例で認識できたと思うが，まだ「推論規則はこれで十分か？」という問いには答えられていない．
	この問いへの解答は，\(\symcal{R}\)における導出がもつ一般的な性質を観察したり他の定式化との等価性を議論するよりほかはない．
	とはいえ現代では本書の定式化と他の数多くの定式化との等価性はすでに示されているため，
	これらの定式化を数学における「計算」の定式化と考えようというのが冒頭で述べたChurchのテーゼである．
}%
．

以上の準備のもと，導出可能性を用いて計算可能関数を定義しよう．

\index[widx]{けいさんかのうかんすう@計算可能関数}
\begin{Def} \label{Def:computablefunction}
	\(n\)を正の整数とする．関数\(f \colon \NaturalNumbers^n \to \NaturalNumbers\)が\term{計算可能関数}であるとは，
	評価式の有限集合\(E\)と関数記号\(\obj{f}\)が存在して，任意の\(x_1, x_2, \dots, x_n, y \in \NaturalNumbers\)に対して
	\begin{equation}
		\apply{f}{x_1, x_2, \dots, x_n} = y
		\quad\metaequivalent\quad
		E \evaluatable \apply{\obj{f}}{\numeral{x_1}, \numeral{x_2}, \dots, \numeral{x_n}} \evaluate \numeral{y}
		\label{eq:computablefunction}
	\end{equation}
	が成り立つことをいう．
\end{Def}

\Cref{eq:computablefunction}の\(\metaequivalent\)は2つの主張を同時に要求していることに注意しよう．
1つは，正しい値\(\apply{f}{x_1, \dots, x_n}\)に対応する評価式\(\apply{\obj{f}}{\numeral{x_1}, \dots, \numeral{x_n}} \evaluate \numeral{\apply{f}{x_1, \dots, x_n}}\)が実際に\(E\)から導出できるという主張であり，
もう1つは，\(\apply{\obj{f}}{\numeral{x_1}, \dots, \numeral{x_n}}\)を左辺とする評価式として\(E\)から導出できる数項はこれ以外に存在しないという主張である．
すなわち，\(E\)は関数記号\(\obj{f}\)を通じて，\(f\)の値を過不足なく計算する手続きを与えていなければならない．

\Cref{Ex:computablefunctionevaluatable}, \cref{Ex:computablefunctionevaluatable2}はまさにこの意味において，
\(\apply{\obj{f}}{a_0, \numeral{0}} \evaluate a_0\)と\(\apply{\obj{f}}{a_0, \apply{\succfunc}{a_1}} \evaluate \apply{\succfunc}{\apply{\obj{f}}{a_0, a_1}}\)からなる\(E\)によって，
加法に相当する関数が計算可能関数であることを（後半の一意性の主張は除いて）具体的に示す例になっている．
この一意性の主張も含めた完全な確認は\cref{Que:additioncomputable}で行う．

以下では，帰納的関数がすべて計算可能関数であることを示す．
そのために，\(E\)からの導出のしかたそのものに関する次の事実を用いる．

\begin{Lemma} \label[Lemma]{lemma:evaluationderivationshape}
	\(E\)を評価式の有限集合とし，評価式\(t_1 \evaluate t_2\)が\(E\)から導出可能であるとする．
	このとき，\(E\)のある要素に(SUB)を有限回（0回でもよい）適用して得られる評価式\(t_1 \evaluate t_2'\)が存在して，
	\(t_2\)は\(t_2'\)に(REPL)を有限回（0回でもよい）適用して得られる．
\end{Lemma}

\begin{proof}
	\(t_1 \evaluate t_2\)を下段とする\(E\)からの導出図の構成（\cref{Def:evaluatable}）に関する帰納法による．
	\(t_1 \evaluate t_2 \in E\)である場合は，これ自身に(SUB), (REPL)をともに0回適用したものと考えればよい．

	(SUB)によって\(t_1 \evaluate t_2\)が得られた場合，すなわちある評価式\(u_1 \evaluate u_2\)を下段とする導出図から\(t_1 \evaluate t_2 = \SubEvaluate{u_1 \evaluate u_2}{a}{n}\)が得られた場合を考える．
	帰納法の仮定より，\(E\)のある要素に(SUB)を有限回適用して得られる評価式\(u_1 \evaluate v\)が存在して，\(u_2\)は\(v\)に(REPL)を有限回適用して得られる．
	(REPL)によって置き換えられる部分項はもとより数項のみからなり変数記号\(a\)を含まないから，
	その置き換えは，さらに変数記号\(a\)を数項\(n\)に置き換える操作の影響を受けない．
	よって，\(v\)にこれらの(REPL)を適用してから\(a\)を\(n\)に置き換えても，先に\(a\)を\(n\)に置き換えてから同じ(REPL)を適用しても，得られる評価式は変わらない．
	そこで，\(u_1 \evaluate v\)に，変数記号\(a\)を数項\(n\)に置き換える(SUB)をさらに適用して得られる評価式を\(t_1 \evaluate t_2'\)とすれば，
	これは\(E\)のある要素に(SUB)を有限回適用して得られたものであり，\(t_2\)はこの\(t_2'\)に，先と同じ(REPL)の適用を行った結果に等しい．

	(REPL)によって\(t_1 \evaluate t_2\)が得られた場合，すなわち左側の前提\(\symcal{D}_1\)の下段が\(t_1 \evaluate w\)であるとする．
	帰納法の仮定を\(\symcal{D}_1\)に適用すれば，\(E\)のある要素に(SUB)を有限回適用して得られる評価式\(t_1 \evaluate t_2'\)が存在して，\(w\)は\(t_2'\)に(REPL)を有限回適用して得られる．
	\(t_1 \evaluate t_2\)は\(t_1 \evaluate w\)にさらに(REPL)を1回適用したものだから，
	\(t_2\)も同じ\(t_2'\)に(REPL)を（1回多く）適用した結果として得られる．
\end{proof}

\Cref{lemma:evaluationderivationshape}からもわかるように，
今回の体系\(\symcal{R}\)は推論規則が非常に単純であるがゆえに極めて限定された形の導出で導出可能な評価式を網羅できる．
これは，導出可能な評価式についても同様である．

\index[widx]{しゅようかんすうきごう@主要関数記号}
\begin{Def} \label{Def:computablefunctionheadfunctionsymbol}
	項\(\apply{f}{t_1, t_2, \dots, t_n}\)（\(f\)は関数記号，\(t_1, \dots, t_n\)は項）に対し，
	関数記号\(f\)をこの項の\term{主要関数記号}と呼ぶ．
	なお，変数記号や定数記号\(\obj{0}\)は主要関数記号を持たない．
\end{Def}

\begin{Lemma} \label[Lemma]{lemma:evaluationreplrestriction}
	\(E\)を評価式の有限集合とし，\(n_1, n_2, \dots, n_r, n\)は数項，\(f\)は関数記号とする．
	\(E\)から\(\apply{f}{n_1, n_2, \dots, n_r} \evaluate n\)の形の評価式が導出可能であるならば，
	\(f\)は\(E\)のある要素の左辺の主要関数記号である．
\end{Lemma}

\begin{proof}
	\(E \evaluatable \apply{k}{n_1, \dots, n_r} \evaluate n\)であるとする．
	\cref{lemma:evaluationderivationshape}より，\(E\)のある要素に(SUB)を有限回適用すると，その左辺として\(\apply{k}{n_1, \dots, n_r}\)が得られる．
	(SUB)は変数記号を数項に置き換えるだけの操作であり，項の主要関数記号を変えることはないから（\cref{Def:computablefunctionheadfunctionsymbol}），
	この\(E\)の要素の左辺の主要関数記号は\(\apply{k}{n_1, \dots, n_r}\)の主要関数記号，すなわち\(k\)と一致する．
\end{proof}

\Cref{lemma:evaluationreplrestriction}の対偶をとれば，関数記号\(k\)がどの要素の左辺の主要関数記号でもないならば，
\(E\)から\(\apply{k}{n_1, \dots, n_r} \evaluate n\)の形の評価式は導出できない．
すなわち，(REPL)によって数項に置き換えられる部分項\(\apply{k}{n_1, \dots, n_r}\)の主要関数記号\(k\)は，
つねに\(E\)のいずれかの要素の左辺の主要関数記号として現れるものに限られる．

\begin{Que} \label[Que]{Que:basicfunctionscomputable}
	\cref{lemma:evaluationderivationshape}, \cref{lemma:evaluationreplrestriction}を用いて，
	ゼロ関数\(\zero^n\)，後者関数\(\suc\)，射影関数\(\proj_i^n\)（\cref{Def:recursivefunction}の規則1--3）が，
	いずれも計算可能関数であることを示せ．
\end{Que}

\begin{Que} \label[Que]{Que:additioncomputable}
	\cref{Ex:computablefunctionevaluatable}, \cref{Ex:computablefunctionevaluatable2}における評価式の集合\(E\)と関数記号\(\obj{f}\)を考える．
	\cref{lemma:evaluationderivationshape}, \cref{lemma:evaluationreplrestriction}を用いて，任意の\(x, y, z \in \NaturalNumbers\)に対して
	\[
		x + y = z
		\quad\metaequivalent\quad
		E \evaluatable \apply{\obj{f}}{\numeral{x}, \numeral{y}} \evaluate \numeral{z}
	\]
	が成り立つことを示せ．
	すなわち，加法を表す関数\(\pair{x, y} \mapsto x + y\)が，この\(E\)と\(\obj{f}\)によって\cref{Def:computablefunction}の意味で計算可能関数であることを確かめよ．
\end{Que}

\begin{Que} \label[Que]{Que:fibonaccicomputable}
	自然数の列\(\apply{\Fib}{0}, \apply{\Fib}{1}, \apply{\Fib}{2}, \dots\)（フィボナッチ数列）を
	\begin{align*}
		\apply{\Fib}{0}     & = 0,                                                                             \\
		\apply{\Fib}{1}     & = 1,                                                                             \\
		\apply{\Fib}{n + 2} & = \apply{\Fib}{n + 1} + \apply{\Fib}{n} \quad \text{(\(n \in \NaturalNumbers\))}
	\end{align*}
	として定める．\(\Fib \colon \NaturalNumbers \to \NaturalNumbers\)が計算可能関数であることを示せ．

	\(\apply{\Fib}{n}\)だけを1変数の原始帰納法の形（\cref{Def:recursivefunction}規則5）に直接乗せようとすると，
	1つ前の値だけでなく2つ前の値も必要になるためうまくいかない．
	そこで，組\(\pair{\apply{\Fib}{n}, \apply{\Fib}{n+1}}\)を求める2つの関数記号を，互いに参照し合う形で定義することを考えよ．
	その際，加法を表す関数記号として\cref{Que:additioncomputable}の\(\obj{f}\)（とその\(E\)）を部品として使ってよい．
\end{Que}

帰納的関数がすべて計算可能関数であることを，2段階に分けて示そう．
まずは，\cref{Def:recursivefunction}の規則6（\(\muoperator\)作用素）を用いない場合，すなわち原始帰納的関数の場合を考える．

\begin{Thm} \label[Thm]{Thm:primitiverecursivefunctioncomputable}
	任意の原始帰納的関数は計算可能関数である．
\end{Thm}

\begin{proof}
	原始帰納的関数の構成に関する帰納法，すなわち\cref{Def:recursivefunction}の規則1--5に関する帰納法によって示す．
	すなわち，これらの規則のそれぞれについて，用いられている関数がすべて計算可能関数であるならば，その規則によって得られる関数も計算可能関数であることを示せばよい．
	規則1（ゼロ関数），規則2（後者関数），規則3（射影関数）については\cref{Que:basicfunctionscomputable}ですでに示した．
	なお，関数記号は可算無限に存在するため，以下ではそれまでに登場したいずれの評価式の集合にも現れない新しい関数記号を自由に選べることに注意する．

	規則4（合成）について．\(g_1, \dots, g_m \colon \NaturalNumbers^n \to \NaturalNumbers\)と\(h \colon \NaturalNumbers^m \to \NaturalNumbers\)が計算可能関数であるとし，
	帰納法の仮定によりそれぞれに対応する評価式の集合\(E_{g_1}, \dots, E_{g_m}, E_h\)と互いに相異なる関数記号\(\obj{g_1}, \dots, \obj{g_m}, \obj{h}\)がとれるとする．
	\(a_1, \dots, a_n\)を相異なる変数記号とし，新しい関数記号\(\obj{f}\)を用いて
	\[
		E = E_{g_1} \cup \dots \cup E_{g_m} \cup E_h \cup
		\Set{\apply{\obj{f}}{a_1, \dots, a_n} \evaluate \apply{\obj{h}}{\apply{\obj{g_1}}{a_1, \dots, a_n}, \dots, \apply{\obj{g_m}}{a_1, \dots, a_n}}}
	\]
	とする．\(x_1, \dots, x_n\)を固定する．
	\(\obj{f}\)は新しい関数記号としたので，\cref{lemma:evaluationderivationshape}より，
	\(E \evaluatable \apply{\obj{f}}{\numeral{x_1}, \dots, \numeral{x_n}} \evaluate n\)となる数項\(n\)は，
	\(E\)の最後の要素に(SUB)を\(n\)回適用して得られる評価式
	\[
		\apply{\obj{f}}{\numeral{x_1}, \dots, \numeral{x_n}} \evaluate
		\apply{\obj{h}}{\apply{\obj{g_1}}{\numeral{x_1}, \dots, \numeral{x_n}}, \dots, \apply{\obj{g_m}}{\numeral{x_1}, \dots, \numeral{x_n}}}
	\]
	の右辺に，(REPL)を有限回適用して得られるものに限られる．
	\cref{lemma:evaluationreplrestriction}より，この(REPL)によって置き換えられる部分項の主要関数記号は\(\obj{g_1}, \dots, \obj{g_m}, \obj{h}\)のいずれかであり，
	帰納法の仮定（各\(g_i, h\)についての\cref{eq:computablefunction}）により，
	部分項\(\apply{\obj{g_i}}{\numeral{x_1}, \dots, \numeral{x_n}}\)を置き換えられる数項は\(\numeral{\apply{g_i}{x_1, \dots, x_n}}\)しかなく，
	これらがすべて置き換えられたのちに残る部分項を置き換えられる数項も
	\[
		\numeral{\apply{h}{\apply{g_1}{x_1, \dots, x_n}, \dots, \apply{g_m}{x_1, \dots, x_n}}} = \numeral{\apply{f}{x_1, \dots, x_n}}
	\]
	しかない．よって\(n = \numeral{\apply{f}{x_1, \dots, x_n}}\)である．
	逆に，帰納法の仮定より\(E \evaluatable \apply{\obj{g_i}}{\numeral{x_1}, \dots, \numeral{x_n}} \evaluate \numeral{\apply{g_i}{x_1, \dots, x_n}}\)（各\(i\)）と
	\[
		E \evaluatable \apply{\obj{h}}{\numeral{\apply{g_1}{x_1, \dots, x_n}}, \dots, \numeral{\apply{g_m}{x_1, \dots, x_n}}} \evaluate \numeral{\apply{f}{x_1, \dots, x_n}}
	\]
	が成り立つから，これらを上記の(SUB)の適用結果と(REPL)によって\(m + 1\)回組み合わせれば，
	実際に\(E \evaluatable \apply{\obj{f}}{\numeral{x_1}, \dots, \numeral{x_n}} \evaluate \numeral{\apply{f}{x_1, \dots, x_n}}\)が導出できる．

	規則5（原始帰納法）について．\(g \colon \NaturalNumbers^n \to \NaturalNumbers\)と\(h \colon \NaturalNumbers^{n+2} \to \NaturalNumbers\)が計算可能関数であるとし，
	帰納法の仮定により対応する評価式の集合\(E_g, E_h\)と相異なる関数記号\(\obj{g}, \obj{h}\)がとれるとする．
	\(a_1, \dots, a_n, b\)を相異なる変数記号とし，新しい関数記号\(\obj{f}\)を用いて
	\begin{align*}
		E = E_g \cup E_h \cup \Set{
			\begin{aligned}
				 & \apply{\obj{f}}{a_1, \dots, a_n, \numeral{0}} \evaluate \apply{\obj{g}}{a_1, \dots, a_n},                                                 \\
				 & \apply{\obj{f}}{a_1, \dots, a_n, \apply{\succfunc}{b}} \evaluate \apply{\obj{h}}{a_1, \dots, a_n, b, \apply{\obj{f}}{a_1, \dots, a_n, b}}
			\end{aligned}
		}
	\end{align*}
	とする．これは，\cref{Ex:computablefunctionevaluatable}, \cref{Ex:computablefunctionevaluatable2}の\(E\)を，一般の計算可能関数\(\obj{g}, \obj{h}\)に対して一般化したものにほかならない．
	\(x_1, \dots, x_n\)を固定し，\(y\)に関する帰納法によって，
	\(E \evaluatable \apply{\obj{f}}{\numeral{x_1}, \dots, \numeral{x_n}, \numeral{y}} \evaluate n\)となる数項\(n\)がちょうど\(\numeral{\apply{f}{x_1, \dots, x_n, y}}\)であることを示す．

	\(y = 0\)のとき．\(\numeral{0}\)と\(\apply{\succfunc}{\dots}\)の形は互いに排他的である（\cref{Def:computablefunctionnumeral}）から，
	\(\obj{f}\)が新しい関数記号であることと\cref{lemma:evaluationderivationshape}より，\(n\)は\(E\)の1つ目の等式に(SUB)を適用して得られる評価式
	\(\apply{\obj{f}}{\numeral{x_1}, \dots, \numeral{x_n}, \numeral{0}} \evaluate \apply{\obj{g}}{\numeral{x_1}, \dots, \numeral{x_n}}\)
	の右辺への(REPL)の適用の結果に限られる．
	\cref{lemma:evaluationreplrestriction}と帰納法の仮定（\(g\)についての\cref{eq:computablefunction}）より，この右辺を置き換えられる数項は\(\numeral{\apply{g}{x_1, \dots, x_n}} = \numeral{\apply{f}{x_1, \dots, x_n, 0}}\)しかないから，\(n\)はこれに等しい．
	逆に，帰納法の仮定による\(E \evaluatable \apply{\obj{g}}{\numeral{x_1}, \dots, \numeral{x_n}} \evaluate \numeral{\apply{f}{x_1, \dots, x_n, 0}}\)と上記の(SUB)の適用を(REPL)で組み合わせれば，実際にこの評価式は導出できる．

	\(y\)から\(y + 1\)への場合．同様に，\(n\)は\(E\)の2つ目の等式に(SUB)を適用して得られる評価式
	\(\apply{\obj{f}}{\numeral{x_1}, \dots, \numeral{x_n}, \numeral{y+1}} \evaluate \apply{\obj{h}}{\numeral{x_1}, \dots, \numeral{x_n}, \numeral{y}, \apply{\obj{f}}{\numeral{x_1}, \dots, \numeral{x_n}, \numeral{y}}}\)
	の右辺への(REPL)の適用の結果に限られる．
	\cref{lemma:evaluationreplrestriction}より，この右辺で置き換えの対象となりうる部分項は\(\obj{f}\)を主要関数記号とするものと\(\obj{h}\)を主要関数記号とするものだけである．
	まず内側の部分項\(\apply{\obj{f}}{\numeral{x_1}, \dots, \numeral{x_n}, \numeral{y}}\)は，\(y\)についての帰納法の仮定によって\(\numeral{\apply{f}{x_1, \dots, x_n, y}}\)にしか置き換えられない．
	これにより残る部分項\(\apply{\obj{h}}{\numeral{x_1}, \dots, \numeral{x_n}, \numeral{y}, \numeral{\apply{f}{x_1, \dots, x_n, y}}}\)は，帰納法の仮定（\(h\)についての\cref{eq:computablefunction}）によって\(\numeral{\apply{h}{x_1, \dots, x_n, y, \apply{f}{x_1, \dots, x_n, y}}} = \numeral{\apply{f}{x_1, \dots, x_n, y+1}}\)にしか置き換えられない．
	よって\(n = \numeral{\apply{f}{x_1, \dots, x_n, y+1}}\)である．
	逆に，\(y\)についての帰納法の仮定と\(h\)についての帰納法の仮定を上記の(SUB)の適用結果に(REPL)で2回組み合わせれば，実際にこの評価式は導出できる．
	以上の\(y\)に関する帰納法により，すべての\(y\)について主張が成り立つ．

	以上，規則1--5のいずれについても，用いられている関数がすべて計算可能関数であれば，得られる関数も計算可能関数であることが示された．
	\cref{Def:recursivefunction}の規則6を1度も適用せずに得られる帰納的関数が原始帰納的関数であったから，
	原始帰納的関数の構成に関する帰納法によって，任意の原始帰納的関数が計算可能関数であることが従う．
\end{proof}

\Cref{Thm:primitiverecursivefunctioncomputable}の証明のうち規則1--5に関する部分は，
そこで用いた関数が計算可能関数であるという仮定のみを使っており，それらが原始帰納的関数であることは一切用いていない．
したがって，同じ議論は一般の帰納的関数の構成に関する帰納法においても，規則1--5についてそのまま通用する．
あとは，規則6（\(\muoperator\)作用素）についても同様のことが成り立つことを示せば，帰納的関数がすべて計算可能関数であることがわかる．

\begin{Thm} \label[Thm]{Thm:recursivefunctioncomputable}
	任意の帰納的関数は計算可能関数である．
\end{Thm}

\begin{proof}
	帰納的関数の構成に関する帰納法によって示す．
	規則1--5については\cref{Thm:primitiverecursivefunctioncomputable}の証明における議論がそのまま適用できるので，
	以下では規則6についてのみ示せばよい．
	関数記号は可算無限に存在するため，以下ではそれまでに登場したいずれの評価式の集合にも現れない新しい関数記号を自由に選べることに注意する．

	規則6（\(\muoperator\)作用素）について．\(g \colon \NaturalNumbers^{n+1} \to \NaturalNumbers\)が計算可能関数であり，
	すべての\(x_1, \dots, x_n\)に対して\(\apply{g}{x_1, \dots, x_n, y} = 0\)となる\(y\)が存在するとする．
	帰納法の仮定により対応する評価式の集合\(E_g\)と関数記号\(\obj{g}\)がとれる．
	\(a_1, \dots, a_n, b, c\)を相異なる変数記号とし，新しい関数記号\(\obj{f}, \obj{R}, \obj{S}\)を用いて
	\begin{align*}
		E = E_g \cup \Set{
			\begin{aligned}
				 & \apply{\obj{f}}{a_1, \dots, a_n} \evaluate \apply{\obj{R}}{a_1, \dots, a_n, \numeral{0}},                                  \\
				 & \apply{\obj{R}}{a_1, \dots, a_n, b} \evaluate \apply{\obj{S}}{\apply{\obj{g}}{a_1, \dots, a_n, b}, a_1, \dots, a_n, b},    \\
				 & \apply{\obj{S}}{\numeral{0}, a_1, \dots, a_n, b} \evaluate b,                                                              \\
				 & \apply{\obj{S}}{\apply{\succfunc}{c}, a_1, \dots, a_n, b} \evaluate \apply{\obj{R}}{a_1, \dots, a_n, \apply{\succfunc}{b}}
			\end{aligned}
		}
	\end{align*}
	とする．\(\obj{R}\)は「\(b\)以上で最小の，\(g\)が0となる引数」を求める関数記号，\(\obj{S}\)はその探索の1歩分を，\(g\)の値によって場合分けして表す関数記号である．
	\(x_1, \dots, x_n\)を固定し，\(y^{\ast} = \apply{\muoperator y}{\apply{g}{x_1, \dots, x_n, y} = 0}\)（\(= \apply{f}{x_1, \dots, x_n}\)）とおく．

	\(b = y^{\ast}, y^{\ast} - 1, \dots, 0\)の順に，\(y^{\ast} - b\)に関する帰納法によって，
	\(E \evaluatable \apply{\obj{R}}{\numeral{x_1}, \dots, \numeral{x_n}, \numeral{b}} \evaluate n\)となる数項\(n\)がちょうど\(\numeral{y^{\ast}}\)であることを示す．
	\(\obj{R}\)についての等式は\(E\)にただ1つしかないから，\cref{lemma:evaluationderivationshape}より，
	この\(n\)は\(\obj{R}\)の等式に(SUB)を適用して得られる評価式\(\apply{\obj{R}}{\numeral{x_1}, \dots, \numeral{x_n}, \numeral{b}} \evaluate \apply{\obj{S}}{\apply{\obj{g}}{\numeral{x_1}, \dots, \numeral{x_n}, \numeral{b}}, \numeral{x_1}, \dots, \numeral{x_n}, \numeral{b}}\)の右辺への(REPL)の適用の結果に限られる．
	帰納法の仮定（\(g\)についての\cref{eq:computablefunction}）より，部分項\(\apply{\obj{g}}{\numeral{x_1}, \dots, \numeral{x_n}, \numeral{b}}\)は\(\numeral{\apply{g}{x_1, \dots, x_n, b}}\)にしか置き換えられない．

	\(b = y^{\ast}\)のとき，\(\apply{g}{x_1, \dots, x_n, y^{\ast}} = 0\)であるから，残る部分項は\(\apply{\obj{S}}{\numeral{0}, \numeral{x_1}, \dots, \numeral{x_n}, \numeral{y^{\ast}}}\)である．
	\(\numeral{0}\)と\(\apply{\succfunc}{\dots}\)の形は互いに排他的であるから，これに整合する\(\obj{S}\)の等式は1つ目のものに限られ，(SUB)によって\(n = \numeral{y^{\ast}}\)と定まる（右辺は変数記号\(b\)そのものであり，これ以上(REPL)は適用できない）．
	逆に，帰納法の仮定による\(E \evaluatable \apply{\obj{g}}{\numeral{x_1}, \dots, \numeral{x_n}, \numeral{y^{\ast}}} \evaluate \numeral{0}\)と，\(\obj{R}, \obj{S}\)の等式への(SUB)の適用とを(REPL)で組み合わせれば，実際に\(E \evaluatable \apply{\obj{R}}{\numeral{x_1}, \dots, \numeral{x_n}, \numeral{y^{\ast}}} \evaluate \numeral{y^{\ast}}\)が導出できる．

	\(b < y^{\ast}\)のとき，\(y^{\ast}\)の最小性から\(\apply{g}{x_1, \dots, x_n, b} \neq 0\)であり，これをある自然数\(c\)を用いて\(\apply{g}{x_1, \dots, x_n, b} = c + 1\)とかけば，
	残る部分項は\(\apply{\obj{S}}{\apply{\succfunc}{\numeral{c}}, \numeral{x_1}, \dots, \numeral{x_n}, \numeral{b}}\)である．
	今度は2つ目の\(\obj{S}\)の等式に整合するから，(SUB)によってこれは\(\apply{\obj{R}}{\numeral{x_1}, \dots, \numeral{x_n}, \numeral{b + 1}}\)への(REPL)の適用の結果に限られる．
	\(b + 1\)についての帰納法の仮定（\(y^{\ast} - (b+1) < y^{\ast} - b\)）より，この部分項は\(\numeral{y^{\ast}}\)にしか置き換えられないから，\(n = \numeral{y^{\ast}}\)である．
	逆に，\(g(x_1, \dots, x_n, b) = c+1\)の導出可能性と，\(b+1\)についての帰納法の仮定による\(E \evaluatable \apply{\obj{R}}{\numeral{x_1}, \dots, \numeral{x_n}, \numeral{b+1}} \evaluate \numeral{y^{\ast}}\)とを，
	\(\obj{R}, \obj{S}\)の等式への(SUB)の適用結果に(REPL)で2回組み合わせれば，実際に\(E \evaluatable \apply{\obj{R}}{\numeral{x_1}, \dots, \numeral{x_n}, \numeral{b}} \evaluate \numeral{y^{\ast}}\)が導出できる．

	以上の帰納法によって\(b = 0\)の場合を得れば，\(E \evaluatable \apply{\obj{R}}{\numeral{x_1}, \dots, \numeral{x_n}, \numeral{0}} \evaluate n\)となる\(n\)は\(\numeral{y^{\ast}}\)に限られ，かつ実際に導出できる．
	これを\(\obj{f}\)の等式と組み合わせれば，同様の議論によって\(E \evaluatable \apply{\obj{f}}{\numeral{x_1}, \dots, \numeral{x_n}} \evaluate n\)となる数項\(n\)は\(\numeral{y^{\ast}} = \numeral{\apply{f}{x_1, \dots, x_n}}\)に限られ，かつ実際にこれが導出できることがわかる．

	以上により，規則6についても，用いられている関数が計算可能関数であれば，得られる関数も計算可能関数であることが示された．
	これと\cref{Thm:primitiverecursivefunctioncomputable}の証明における規則1--5についての議論を合わせれば，
	\cref{Def:recursivefunction}の規則1--6のいずれについても，用いられている関数がすべて計算可能関数であれば，得られる関数も計算可能関数であることがわかる．
	\cref{Def:recursivefunction}の規則7により帰納的関数は上記の規則を有限回適用して得られる関数に限られるから，
	帰納的関数の構成に関する帰納法によって，任意の帰納的関数が計算可能関数であることが従う．
\end{proof}

\Cref{Thm:recursivefunctioncomputable}の逆，すなわちすべての計算可能関数が帰納的関数であることも実は成り立つ．
しかし，その証明には本書の範囲を超える準備が必要となるため，ここでは扱わない．

\section{有限列のコード化} \label[section]{sec:codizeforfinitesequence}

本節では，のちの章への準備としての自然数からなる有限列のコード化について議論する．
有限列のコード化とは，有限列に対して一意な自然数を割り当てることである．
コード化の手法は異なったものが複数知られているが，本書ではCantorの対関数を使ったものを紹介する．
他の手法として有名なものに素因数分解の一意性を用いたものが挙げられる．

なお，本節の内容が必要になるのは\cref{chap:peanoarithmetic}以降である．
そのため「まずは完全性定理まで最速で入門したい」という場合には本節の内容を読み飛ばしても問題ない．

\index[widx]{Cantorのついかんすう@Cantorの対関数}
\index[sidx]{\(\pairfunc\)：Cantorの対関数}
\begin{Def} \label{Def:cantorpairfunction}
	関数\(\pairfunc \colon \NaturalNumbers^2 \to \NaturalNumbers\)を
	\begin{equation}
		\apply{\pairfunc}{x, y} = \sum_{i < x + y + 1} i + x = \frac{1}{2} \paren{x + y} \paren{x + y + 1} + x
		\quad \text{(\(x, y \in \NaturalNumbers\))}
		\label{eq:cantorpairfunction}
	\end{equation}
	と定める．この関数\(\pairfunc\)を\term{Cantorの対関数}と呼ぶ．
\end{Def}

よく知られているように，Cantorの対関数は\(\NaturalNumbers^2\)から\(\NaturalNumbers\)への全単射になっている．

\begin{Lemma} \label[Lemma]{lemma:cantorpairfunctionbijection}
	Cantorの対関数\(\pairfunc \colon \NaturalNumbers^2 \to \NaturalNumbers\)は全単射であり，
	その逆関数\(x \mapsto \pair{\apply{\pairfunc_1}{x}, \apply{\pairfunc_2}{x}}\)は以下で与えられる：
	\begin{align}
		\apply{t}{x}           & = \max \Set{n \in \NaturalNumbers | \sum_{i < n + 1} i < x}
		\quad \text{(\(x \in \NaturalNumbers\))},
		\label{eq:cantorpairfunctioninversebase}                                             \\
		\apply{\pairfunc_2}{x} & = x - \sum_{i < \apply{t}{x} + 1} i
		\quad \text{(\(x \in \NaturalNumbers\))},
		\label{eq:cantorpairfunctioninverse2}                                                \\
		\apply{\pairfunc_1}{x} & = \apply{t}{x} - \apply{\pairfunc_2}{x}
		\quad \text{(\(x \in \NaturalNumbers\))}.
		\label{eq:cantorpairfunctioninverse1}
	\end{align}
\end{Lemma}

\begin{Def} \label{Def:cantorpairfunctionn}
	正の整数\(n\)に対し，写像\(\pairfunc^n \colon \NaturalNumbers^n \to \NaturalNumbers\)を以下のように帰納的に定義する：
	\begin{align}
		\begin{aligned}
			\apply{\pairfunc^1}{x}                                      & = x
			\quad \text{(\(x \in \NaturalNumbers\))},                                                                                             \\
			\apply{\pairfunc^{n + 1}}{x_0, x_1, \dots, x_{n  - 1}, x_n} & = \apply{\pairfunc}{x_0, \apply{\pairfunc^n}{x_1, \dots, x_{n-1}, x_n}} \\
			                                                            & \quad \text{(\(x_0, x_1, \dots, x_{n-1}, x_n \in \NaturalNumbers\))}.
		\end{aligned}
		\label{eq:cantorpairfunctionn}
	\end{align}
\end{Def}

\begin{Lemma} \label[Lemma]{lemma:cantorpairfunctionbijectionn}
	正の整数\(n\)に対し，\(\pairfunc^n \colon \NaturalNumbers^n \to \NaturalNumbers\)は全単射であり，その逆写像
	\[
		x \mapsto \pair{\apply{\pairfunc_0^n}{x}, \apply{\pairfunc_1^n}{x}, \dots, \apply{\pairfunc_{n-1}^n}{x}}
	\]
	は，\(\pairfunc_2\)を\(i\)回合成して得られる関数\(\pairfunc_2^{\paren{i}}\)を用いて以下のように与えられる：
	\begin{equation}
		\apply{\pairfunc_i^n}{x} =
		\begin{cases*}
			\apply{\pairfunc_1}{x}                                  & (\(i = 0\)のとき，\(x \in \NaturalNumbers\)),         \\
			\apply{\pairfunc_1}{\apply{\pairfunc_2^{\paren{i}}}{x}} & (\(0 < i < n - 1\)のとき，\(x \in \NaturalNumbers\)), \\
			\apply{\pairfunc_2^{\paren{n - 1}}}{x}                  & (\(i = n - 1\)のとき，\(x \in \NaturalNumbers\))
		\end{cases*}
		\label{eq:cantorpairfunctionninv}
	\end{equation}
\end{Lemma}

以上の準備により，有限列のコード化を定義することができる．

\index[sidx]{\(\finitesequencecode{x_0, x_1, \dots, x_{n-1}}\)：有限列のコード}
\index[widx]{こーど@（自然数からなる有限列の）コード}
\begin{Def} \label{Def:codizeforfinitesequence}
	与えられた自然数からなる有限列
	\[
		x_0, x_1, \dots, x_{n-1} \in \bigcup_{n \in \NaturalNumbers} \NaturalNumbers^{\Set{0, 1, \dots, n-1}}
	\]
	に対し，そのコード
	\begin{equation}
		\finitesequencecode{x_0, x_1, \dots, x_{n-1}}
		\label{eq:finitesequencecode}
	\end{equation}
	を%
	\footnote{%
		順序対と記号が被っているが，文脈上混同する可能性は低いと考えられる．
	}%
	\begin{equation}
		\finitesequencecode{x_0, x_1, \dots, x_{n-1}} =
		\begin{cases*}
			\apply{\pairfunc}{0,0}                                                  & (\(n = 0\)のとき), \\
			\apply{\pairfunc}{0, 1 + x_0}                                           & (\(n = 1\)のとき), \\
			\apply{\pairfunc}{n - 1, \apply{\pairfunc^n}{x_0, x_1, \dots, x_{n-1}}} & (それ以外のとき)
		\end{cases*}
		\label{eq:finitesequencecodedefinition}
	\end{equation}
	と定める．
	与えられた有限列が空列である場合もそのコードが0として定まっていることに注意しよう．
\end{Def}

\Cref{lemma:cantorpairfunctionbijectionn}から，以下の主張が従う．

\begin{Lemma} \label[Lemma]{lemma:codizeforfinitesequencebijection}
	与えられた自然数からなる有限列
	\[
		x_0, x_1, \dots, x_{n-1} \in \bigcup_{n \in \NaturalNumbers} \NaturalNumbers^{\Set{0, 1, \dots, n-1}}
	\]
	に対し，そのコード
	\begin{equation*}
		\finitesequencecode{x_0, x_1, \dots, x_{n-1}}
	\end{equation*}
	を対応させる\(\bigcup_{n \in \NaturalNumbers} \NaturalNumbers^{\Set{0,1,\dots,n-1}}\)から\(\NaturalNumbers\)への関数は全単射である．
\end{Lemma}

\index[sidx]{\(\apply{\lh}{x}\)：コード化された有限列の長さ}
\begin{Def} \label{Def:finitesequencecodelength}
	与えられた自然数\(x\)に対し，\(x\)に対応している自然数からなる有限列
	\(x_0, x_1, \dots, x_{n-1}\)の長さ\(n\)を
	\begin{equation}
		\apply{\lh}{x}
		\label{eq:finitesequencecodelength}
	\end{equation}
	と表す．
	\(\apply{\lh}{x}\)は具体的に
	\begin{equation}
		\apply{\lh}{x} =
		\begin{cases*}
			0                          & (\(x = 0\)のとき), \\
			\apply{\pairfunc_1}{x} + 1 & (それ以外のとき)
		\end{cases*}
		\label{eq:finitesequencecodelengthdefinition}
	\end{equation}
	と表示できる．
\end{Def}

\index[sidx]{\(\codefunctioninv{x}{i}\)：コードから有限列を得る関数}
\begin{Def} \label{Def:codefunctioninv}
	与えられた1以上の自然数\(x\)に対し，\(x\)がそのコードとなる自然数からなる有限列
	\(x_0, x_1, \dots, x_{n-1}\)における\(x_i\)を
	\begin{equation}
		\codefunctioninv{x}{i}
		\label{eq:codefunctioninv}
	\end{equation}
	と表す．
	\(\codefunctioninv{x}{i}\)は具体的に
	\begin{equation}
		\codefunctioninv{x}{i} =
		\begin{cases*}
			\apply{\pairfunc_2}{x} - 1                                   & (\(\apply{\lh}{x} = 1\)のとき) \\
			\apply{\pairfunc_i^{\apply{\lh}{x}}}{\apply{\pairfunc_2}{x}} & (それ以外のとき)
		\end{cases*}
		\label{eq:codefunctioninvdefinition}
	\end{equation}
	と表示できる．

	\(x = 0\)である，もしくは\(\apply{\lh}{x} \leq i\)である場合には\cref{eq:codefunctioninv}の値は0と約束する．
\end{Def}

\index[sidx]{\(x \codeconnection y\)：有限列の連結に対応するコード}
\begin{Def} \label{Def:codeconnection}
	自然数\(x, y\)を有限列のコードとして\(x = \finitesequencecode{x_0, x_1, \dots, x_{n-1}}\)，
	\(y = \finitesequencecode{y_0, y_1, \dots, y_{m-1}}\)
	と表示し，それを用いて
	\begin{equation}
		x \codeconnection y = \finitesequencecode{x_0, x_1, \dots, x_{n-1}, y_0, y_1, \dots, y_{m-1}}
		\label{eq:codeconnection}
	\end{equation}
	と定める．
\end{Def}

\begin{Lemma} \label[Lemma]{lemma:codeconnection}
	関数\(\connectionfunc \colon \NaturalNumbers^3 \to \NaturalNumbers\)を以下のように帰納的に定義する：
	\begin{align}
		\begin{aligned}
			\apply{\connectionfunc}{0, x, y}   & = y \quad \text{(\(x, y \in \NaturalNumbers\))}, \\
			\apply{\connectionfunc}{n+1, x, y} & =
			\begin{cases*}
				\apply{\pairfunc}{\apply{\pairfunc^{\apply{\lh}{x}}_{\apply{\lh}{x}-n-1}}{x},\apply{\connectionfunc}{n, x, y}} & (\(n < \apply{\lh}{x}\)のとき) \\
				\apply{\pairfunc}{0,\apply{\connectionfunc}{n, x, y}}                                                          & (それ以外のとき)                   \\
			\end{cases*}
		\end{aligned}
		\label{eq:codeconnectionfunc}
	\end{align}

	\(x, y \in \NaturalNumbers\)に対する\(x \codeconnection y\)は，上記\(\connectionfunc\)を用いて
	\begin{equation}
		x \codeconnection y = \apply{\connectionfunc}{\apply{\lh}{x}, x, y}
		\label{eq:codeconnectionconnectionfunc}
	\end{equation}
	と表せる．
\end{Lemma}


\Cref{chap:peanoarithmetic}での議論においては，
ここで定義した種々の関数がCantorの対関数を使って具体的に表示できているという点も重要となる．